% =============================================================================
% Section 9: Conclusion
% =============================================================================
\section{Conclusion}\label{sec:conclusion}

This paper has presented the first comprehensive empirical application
of the Import Strategicness Index (ISI) to the~27 member states of
the European Union.  Drawing on a unified data snapshot (ISI~v1.0,
reference year~2024), we have computed composite scores for all
countries, decomposed the composite into its six constituent axes,
and subjected the resulting ranking to a four-part robustness
programme.

\subsection{Summary of contributions}

The paper makes seven distinct contributions:

\begin{enumerate}[leftmargin=2em]
  \item It produces the first complete \textbf{ISI composite ranking}
    for the EU-27, identifying Malta as the most concentrated and
    France as the least concentrated member state.

  \item It provides \textbf{distributional diagnostics} for the
    composite (histogram, ECDF, Lorenz curve with Gini = \num{0.1155}),
    establishing that the EU-27 distribution is compressed and
    unimodal.

  \item It reports \textbf{axis-level statistics}, identifying defence
    and logistics as the axes with the highest dispersion and the
    financial axis as the most compressed.

  \item It documents the \textbf{axis dominance structure}: defence is
    the dominant axis for 16~countries, logistics for 11, and no other
    axis is ever the maximum.

  \item It computes a \textbf{variance contribution decomposition},
    finding that defence (35.10\%) and logistics (34.26\%) jointly
    account for 69.36\% of composite variance.

  \item It executes a \textbf{four-exercise robustness programme}
    (LOO, z-score, Dirichlet MC, winsorization) and finds the
    ranking robust to outlier capping ($\rho = 1.000$) and random
    weight perturbation ($\bar{\rho} = 0.979$), with moderate
    sensitivity to defence-axis removal ($\rho = 0.833$) and
    z-score standardisation ($\rho = 0.896$).

  \item It profiles \textbf{eight representative countries} via
    radar charts, illustrating three archetypal concentration
    patterns: small-island, sector-specific spike, and broad
    diversification.
\end{enumerate}

\subsection{Directions for future research}

Several extensions are planned:

\begin{itemize}[leftmargin=2em]
  \item \textbf{Time-series extension.}  Computing the ISI for
    multiple reference years (2015--2024) to analyse trend dynamics,
    structural breaks, and the effect of geopolitical shocks (e.g.,
    the 2022 Russia--Ukraine conflict) on concentration patterns.

  \item \textbf{Geographic expansion.}  Extending the ISI to G-20
    or OECD countries to enable broader cross-country comparisons
    and identify global concentration patterns.

  \item \textbf{Sub-indicator disaggregation.}  Reporting
    sub-indicator-level results for each axis to provide more
    granular policy guidance.

  \item \textbf{Endogenous weighting.}  Exploring data-driven
    weighting schemes (e.g., principal-component weights, benefit-of-
    the-doubt weights) as alternatives to equal weighting.

  \item \textbf{Interaction with macroeconomic outcomes.}  Testing
    whether high ISI scores are associated with greater GDP
    volatility, trade disruptions, or supply-chain bottlenecks.
\end{itemize}

\subsection{Closing remark}

The ISI is designed as a transparent, reproducible measurement tool
for strategic import dependence.  Its value lies not in producing a
single ``correct'' ranking but in providing a structured empirical
basis for policy deliberation.  The robustness programme demonstrates
that the ranking is stable under most perturbations, and the
axis-level disaggregation enables targeted analysis of specific
vulnerability domains.  We hope that this empirical application will
serve as a reference for researchers and policy-makers working on
European strategic autonomy.
